% !TeX root = ../教材题库_全部条目_转写工作稿.tex
% Source: XJ-B1-C06-PR-P220-Y0475
% 第6章 / 6.2 抽样 / 教材 p.220 / PDF 第227页
% OCR未校对：公式与文字均可能有误；未完成数学TeX转写；禁止直接发给学生
% 原题及相邻原解截图（核对用）：../题目原图/XJ-B1-C06-PR-P220-Y0475-01.webp
% 本文件只在转写完成后作为题干使用；OCR 条目尚不是可靠题干。
\begin{bankitem}{p.220 练习 第2题}{第6章 / 6.2 抽样 / 教材 p.220 / PDF 第227页}{XJ-B1-C06-PR-P220-Y0475}{OCR未校对：公式与文字均可能有误；未完成数学TeX转写；禁止直接发给学生}
\begin{ocrdraft}
[教材p.220 / PDF第227页]
\par
2．某公司甲、乙、丙、丁四个地区分别有150个，120个，180个，150个
\par
销售点．公司为了调查产品销售的情况，需从这600个销售点中抽取一个容量为
\par
100的样本，记这项调查为（1）；在丙地区中有20个特大型销售点，要从中抽取7
\par
个调查其收入和售后服务等情况，记这项调查为（2）．完成（1)（2）这两项调查时，
\par
宜分别采用何种抽样方法？
\end{ocrdraft}
\end{bankitem}
